%% ============================================================
%% Section 7: Conclusions (V0.3)
%% Target: ~400 words
%% New section: synthesis, key takeaways, limitations, call to action
%% ============================================================

\section{Conclusions}
\label{sec:conclusions}

This survey has traced the security and safety landscape of AI-Generated Content from its origins in content-level risks to the emerging frontier of autonomous agentic exploitation. A central meta-pattern emerges: \emph{temporal misalignment}. Attack surfaces evolve through a well-defined progression---direct prompt injection, jailbreaking, indirect prompt injection, and agentic exploitation via tool-invocation ecosystems---while defensive capabilities lag by months and governance frameworks by years. The first agentic AI products shipped in October 2024; the first government-level agentic governance framework appeared in January 2026, a 16-month gap during which AI agents with tool-invocation privileges were deployed at scale under no agent-specific standard~\cite{imda2026agentic}. The ``deploy first, retrofit security later'' pattern, documented throughout Sections~\ref{sec:threats}--\ref{sec:governance}, has extended from content risks to action risks with qualitatively higher consequences.

Key insights from each dimension reinforce this assessment. Generative capabilities have converged across modalities at near-zero marginal cost, with autonomous agency representing a qualitative shift from bounded output to unbounded action (Section~\ref{sec:capabilities}). The threat landscape spans content weaponization, model-level attacks, adversarial manipulation, and societal harms, with the attack surface evolution ladder from prompt injection to agentic exploitation constituting the most consequential and least-governed vector (Section~\ref{sec:threats}). Technical countermeasures face fundamental limits---detection collapses under distribution shift, watermarking confronts information-theoretic impossibility results, alignment is superficial and fragile, and agentic defenses lack compositional guarantees (Section~\ref{sec:countermeasures}). Governance has fractured into three irreconcilable paradigms with zero coverage of agentic AI systems (Section~\ref{sec:governance}). Six open research frontiers remain largely unsolved, with securing autonomous AI agents and unified multi-modal detection identified as the most urgent (Section~\ref{sec:future}).

\paragraph{Limitations.}
This survey has several limitations that readers should consider. First, our literature coverage is predominantly English-language; important contributions in Chinese, Korean, and other languages may be underrepresented. Second, several key data points (e.g., 82.6\% AI-generated phishing~\cite{knowbe4_2025}, 72.8\% MCPTox ASR~\cite{wang2025mcptox}) originate from industry reports rather than peer-reviewed studies, with potential methodological limitations and conflicts of interest. Third, the knowledge cutoff of March 2026 constrains currency in a field where capabilities, threats, and regulations evolve on monthly timescales. Fourth, the agentic AI security literature is extremely young (predominantly 2025--2026), and long-term trend assessments carry inherent uncertainty.

\paragraph{Looking forward.}
The asymmetry between offensive and defensive capabilities demands a paradigm shift from reactive to proactive security. Cross-disciplinary collaboration---spanning computer science, law, ethics, and public policy---is no longer aspirational but operationally necessary. The transition from ``AI generates content'' to ``AI executes actions'' represents both the defining challenge and the most consequential research frontier of the coming decade.
